\IfFileExists{stacks-project.cls}{%
\documentclass{stacks-project}
}{%
\documentclass{amsart}
}

% For dealing with references we use the comment environment
\usepackage{verbatim}
\newenvironment{reference}{\comment}{\endcomment}
%\newenvironment{reference}{}{}
\newenvironment{slogan}{\comment}{\endcomment}
\newenvironment{history}{\comment}{\endcomment}

% For commutative diagrams we use Xy-pic
\usepackage[all]{xy}

% We use 2cell for 2-commutative diagrams.
\xyoption{2cell}
\UseAllTwocells

% We use multicol for the list of chapters between chapters
\usepackage{multicol}

% This is generally recommended for better output
\usepackage{lmodern}
\usepackage[T1]{fontenc}

% For cross-file-references
\usepackage{xr-hyper}

% Package for hypertext links:
\usepackage{hyperref}

% For any local file, say "hello.tex" you want to link to please
% use \externaldocument[hello-]{hello}
\externaldocument[introduction-]{introduction}
\externaldocument[conventions-]{conventions}
\externaldocument[sets-]{sets}
\externaldocument[categories-]{categories}
\externaldocument[topology-]{topology}
\externaldocument[sheaves-]{sheaves}
\externaldocument[sites-]{sites}
\externaldocument[stacks-]{stacks}
\externaldocument[fields-]{fields}
\externaldocument[algebra-]{algebra}
\externaldocument[brauer-]{brauer}
\externaldocument[homology-]{homology}
\externaldocument[derived-]{derived}
\externaldocument[simplicial-]{simplicial}
\externaldocument[more-algebra-]{more-algebra}
\externaldocument[smoothing-]{smoothing}
\externaldocument[modules-]{modules}
\externaldocument[sites-modules-]{sites-modules}
\externaldocument[injectives-]{injectives}
\externaldocument[cohomology-]{cohomology}
\externaldocument[sites-cohomology-]{sites-cohomology}
\externaldocument[dga-]{dga}
\externaldocument[dpa-]{dpa}
\externaldocument[sdga-]{sdga}
\externaldocument[hypercovering-]{hypercovering}
\externaldocument[schemes-]{schemes}
\externaldocument[constructions-]{constructions}
\externaldocument[properties-]{properties}
\externaldocument[morphisms-]{morphisms}
\externaldocument[coherent-]{coherent}
\externaldocument[divisors-]{divisors}
\externaldocument[limits-]{limits}
\externaldocument[varieties-]{varieties}
\externaldocument[topologies-]{topologies}
\externaldocument[descent-]{descent}
\externaldocument[perfect-]{perfect}
\externaldocument[more-morphisms-]{more-morphisms}
\externaldocument[flat-]{flat}
\externaldocument[groupoids-]{groupoids}
\externaldocument[more-groupoids-]{more-groupoids}
\externaldocument[etale-]{etale}
\externaldocument[chow-]{chow}
\externaldocument[intersection-]{intersection}
\externaldocument[pic-]{pic}
\externaldocument[weil-]{weil}
\externaldocument[adequate-]{adequate}
\externaldocument[dualizing-]{dualizing}
\externaldocument[duality-]{duality}
\externaldocument[discriminant-]{discriminant}
\externaldocument[derham-]{derham}
\externaldocument[local-cohomology-]{local-cohomology}
\externaldocument[algebraization-]{algebraization}
\externaldocument[curves-]{curves}
\externaldocument[resolve-]{resolve}
\externaldocument[models-]{models}
\externaldocument[functors-]{functors}
\externaldocument[equiv-]{equiv}
\externaldocument[pione-]{pione}
\externaldocument[etale-cohomology-]{etale-cohomology}
\externaldocument[proetale-]{proetale}
\externaldocument[relative-cycles-]{relative-cycles}
\externaldocument[more-etale-]{more-etale}
\externaldocument[trace-]{trace}
\externaldocument[crystalline-]{crystalline}
\externaldocument[spaces-]{spaces}
\externaldocument[spaces-properties-]{spaces-properties}
\externaldocument[spaces-morphisms-]{spaces-morphisms}
\externaldocument[decent-spaces-]{decent-spaces}
\externaldocument[spaces-cohomology-]{spaces-cohomology}
\externaldocument[spaces-limits-]{spaces-limits}
\externaldocument[spaces-divisors-]{spaces-divisors}
\externaldocument[spaces-over-fields-]{spaces-over-fields}
\externaldocument[spaces-topologies-]{spaces-topologies}
\externaldocument[spaces-descent-]{spaces-descent}
\externaldocument[spaces-perfect-]{spaces-perfect}
\externaldocument[spaces-more-morphisms-]{spaces-more-morphisms}
\externaldocument[spaces-flat-]{spaces-flat}
\externaldocument[spaces-groupoids-]{spaces-groupoids}
\externaldocument[spaces-more-groupoids-]{spaces-more-groupoids}
\externaldocument[bootstrap-]{bootstrap}
\externaldocument[spaces-pushouts-]{spaces-pushouts}
\externaldocument[spaces-chow-]{spaces-chow}
\externaldocument[groupoids-quotients-]{groupoids-quotients}
\externaldocument[spaces-more-cohomology-]{spaces-more-cohomology}
\externaldocument[spaces-simplicial-]{spaces-simplicial}
\externaldocument[spaces-duality-]{spaces-duality}
\externaldocument[formal-spaces-]{formal-spaces}
\externaldocument[restricted-]{restricted}
\externaldocument[spaces-resolve-]{spaces-resolve}
\externaldocument[formal-defos-]{formal-defos}
\externaldocument[defos-]{defos}
\externaldocument[cotangent-]{cotangent}
\externaldocument[examples-defos-]{examples-defos}
\externaldocument[algebraic-]{algebraic}
\externaldocument[examples-stacks-]{examples-stacks}
\externaldocument[stacks-sheaves-]{stacks-sheaves}
\externaldocument[criteria-]{criteria}
\externaldocument[artin-]{artin}
\externaldocument[quot-]{quot}
\externaldocument[stacks-properties-]{stacks-properties}
\externaldocument[stacks-morphisms-]{stacks-morphisms}
\externaldocument[stacks-limits-]{stacks-limits}
\externaldocument[stacks-cohomology-]{stacks-cohomology}
\externaldocument[stacks-perfect-]{stacks-perfect}
\externaldocument[stacks-introduction-]{stacks-introduction}
\externaldocument[stacks-more-morphisms-]{stacks-more-morphisms}
\externaldocument[stacks-geometry-]{stacks-geometry}
\externaldocument[moduli-]{moduli}
\externaldocument[moduli-curves-]{moduli-curves}
\externaldocument[examples-]{examples}
\externaldocument[exercises-]{exercises}
\externaldocument[guide-]{guide}
\externaldocument[desirables-]{desirables}
\externaldocument[coding-]{coding}
\externaldocument[obsolete-]{obsolete}
\externaldocument[fdl-]{fdl}
\externaldocument[index-]{index}

% Theorem environments.
%
\theoremstyle{plain}
\newtheorem{theorem}[subsection]{Theorem}
\newtheorem{proposition}[subsection]{Proposition}
\newtheorem{lemma}[subsection]{Lemma}

\theoremstyle{definition}
\newtheorem{definition}[subsection]{Definition}
\newtheorem{example}[subsection]{Example}
\newtheorem{exercise}[subsection]{Exercise}
\newtheorem{situation}[subsection]{Situation}

\theoremstyle{remark}
\newtheorem{remark}[subsection]{Remark}
\newtheorem{remarks}[subsection]{Remarks}

\numberwithin{equation}{subsection}

% Macros
%
\def\lim{\mathop{\mathrm{lim}}\nolimits}
\def\colim{\mathop{\mathrm{colim}}\nolimits}
\def\Spec{\mathop{\mathrm{Spec}}}
\def\Hom{\mathop{\mathrm{Hom}}\nolimits}
\def\Ext{\mathop{\mathrm{Ext}}\nolimits}
\def\SheafHom{\mathop{\mathcal{H}\!\mathit{om}}\nolimits}
\def\SheafExt{\mathop{\mathcal{E}\!\mathit{xt}}\nolimits}
\def\Sch{\mathit{Sch}}
\def\Mor{\mathop{\mathrm{Mor}}\nolimits}
\def\Ob{\mathop{\mathrm{Ob}}\nolimits}
\def\Sh{\mathop{\mathit{Sh}}\nolimits}
\def\NL{\mathop{N\!L}\nolimits}
\def\CH{\mathop{\mathrm{CH}}\nolimits}
\def\proetale{{pro\text{-}\acute{e}tale}}
\def\etale{{\acute{e}tale}}
\def\QCoh{\mathit{QCoh}}
\def\Ker{\mathop{\mathrm{Ker}}}
\def\Im{\mathop{\mathrm{Im}}}
\def\Coker{\mathop{\mathrm{Coker}}}
\def\Coim{\mathop{\mathrm{Coim}}}

% Boxtimes
%
\DeclareMathSymbol{\boxtimes}{\mathbin}{AMSa}{"02}

%
% Macros for moduli stacks/spaces
%
\def\QCohstack{\mathcal{QC}\!\mathit{oh}}
\def\Cohstack{\mathcal{C}\!\mathit{oh}}
\def\Spacesstack{\mathcal{S}\!\mathit{paces}}
\def\Quotfunctor{\mathrm{Quot}}
\def\Hilbfunctor{\mathrm{Hilb}}
\def\Curvesstack{\mathcal{C}\!\mathit{urves}}
\def\Polarizedstack{\mathcal{P}\!\mathit{olarized}}
\def\Complexesstack{\mathcal{C}\!\mathit{omplexes}}
% \Pic is the operator that assigns to X its picard group, usage \Pic(X)
% \Picardstack_{X/B} denotes the Picard stack of X over B
% \Picardfunctor_{X/B} denotes the Picard functor of X over B
\def\Pic{\mathop{\mathrm{Pic}}\nolimits}
\def\Picardstack{\mathcal{P}\!\mathit{ic}}
\def\Picardfunctor{\mathrm{Pic}}
\def\Deformationcategory{\mathcal{D}\!\mathit{ef}}
\usepackage{cleveref}
\crefname{originaldefinition}{original definition}{original definitions}
\Crefname{originaldefinition}{Original definition}{Original definitions}
\makeatletter
\newlabel{situation-coherent}{{08KB}{source}{}{sourcepointer-1}{}}
\newlabel{situation-coherent@cref}{{[originaldefinition][0][]08KB}{[1][0][]source}{}{}{}}
\makeatother
\crefname{originaldefinition}{original definition}{original definitions}
\Crefname{originaldefinition}{Original definition}{Original definitions}
\makeatletter
\newlabel{artin-definition-openness-versality}{{07XQ}{source}{}{sourcepointer-2}{}}
\newlabel{artin-definition-openness-versality@cref}{{[originaldefinition][0][]07XQ}{[1][0][]source}{}{}{}}
\makeatother
\crefname{originallemma}{original lemma}{original lemmas}
\Crefname{originallemma}{Original lemma}{Original lemmas}
\makeatletter
\newlabel{artin-lemma-get-openness-obstruction-theory}{{0CYF}{source}{}{sourcepointer-3}{}}
\newlabel{artin-lemma-get-openness-obstruction-theory@cref}{{[originallemma][0][]0CYF}{[1][0][]source}{}{}{}}
\makeatother
\crefname{originallemma}{original lemma}{original lemmas}
\Crefname{originallemma}{Original lemma}{Original lemmas}
\makeatletter
\newlabel{lemma-coherent-diagonal}{{08W6}{source}{}{sourcepointer-4}{}}
\newlabel{lemma-coherent-diagonal@cref}{{[originallemma][0][]08W6}{[1][0][]source}{}{}{}}
\makeatother
\crefname{originallemma}{original lemma}{original lemmas}
\Crefname{originallemma}{Original lemma}{Original lemmas}
\makeatletter
\newlabel{lemma-coherent-RS-star}{{08LQ}{source}{}{sourcepointer-5}{}}
\newlabel{lemma-coherent-RS-star@cref}{{[originallemma][0][]08LQ}{[1][0][]source}{}{}{}}
\makeatother
\crefname{originallemma}{original lemma}{original lemmas}
\Crefname{originallemma}{Original lemma}{Original lemmas}
\makeatletter
\newlabel{lemma-coherent-limits}{{08KD}{source}{}{sourcepointer-6}{}}
\newlabel{lemma-coherent-limits@cref}{{[originallemma][0][]08KD}{[1][0][]source}{}{}{}}
\makeatother
\crefname{originaldefinition}{original definition}{original definitions}
\Crefname{originaldefinition}{Original definition}{Original definitions}
\makeatletter
\newlabel{artin-definition-obstruction-theory}{{07YG}{source}{}{sourcepointer-7}{}}
\newlabel{artin-definition-obstruction-theory@cref}{{[originaldefinition][0][]07YG}{[1][0][]source}{}{}{}}
\makeatother
\crefname{originallemma}{original lemma}{original lemmas}
\Crefname{originallemma}{Original lemma}{Original lemmas}
\makeatletter
\newlabel{defos-lemma-flat-ringed-topoi}{{08VW}{source}{}{sourcepointer-8}{}}
\newlabel{defos-lemma-flat-ringed-topoi@cref}{{[originallemma][0][]08VW}{[1][0][]source}{}{}{}}
\makeatother
\crefname{originallemma}{original lemma}{original lemmas}
\Crefname{originallemma}{Original lemma}{Original lemmas}
\makeatletter
\newlabel{defos-lemma-functorial-ringed-topoi}{{0CYE}{source}{}{sourcepointer-9}{}}
\newlabel{defos-lemma-functorial-ringed-topoi@cref}{{[originallemma][0][]0CYE}{[1][0][]source}{}{}{}}
\makeatother
\crefname{originallemma}{original lemma}{original lemmas}
\Crefname{originallemma}{Original lemma}{Original lemmas}
\makeatletter
\newlabel{spaces-perfect-lemma-compute-ext}{{08JR}{source}{}{sourcepointer-10}{}}
\newlabel{spaces-perfect-lemma-compute-ext@cref}{{[originallemma][0][]08JR}{[1][0][]source}{}{}{}}
\makeatother
\crefname{originallemma}{original lemma}{original lemmas}
\Crefname{originallemma}{Original lemma}{Original lemmas}
\makeatletter
\newlabel{more-algebra-lemma-pseudo-coherent-tensor}{{0CYB}{source}{}{sourcepointer-11}{}}
\newlabel{more-algebra-lemma-pseudo-coherent-tensor@cref}{{[originallemma][0][]0CYB}{[1][0][]source}{}{}{}}
\makeatother
\title{Openness of versality for the stack of coherent sheaves in the flat case}
\author{The Stacks Project Authors}
\date{stacks-project:quot.tex:local-snapshot-20261012; GNU FDL 1.2 or later}
\begin{document}
\maketitle
\setcounter{subsection}{1}
\noindent\textit{Editorial note.} All delivered mathematical text is the author's own, copied from the pinned Stacks Project chapter file quot.tex; no statement or proof step was rewritten, reordered or removed. The counted span is the author environment and its own proof environment. References inside the span are handed to the assembler as navigation stubs that carry the canonical tag of the cited tag only; whether each of them is load-bearing is recorded on the stub itself.\par\medskip
\noindent\textit{Allowed standard prerequisites (editorial).}\begin{itemize}
\item Situation situation-coherent (canonical tag 08KB). cited at statement line 1547. Standing situation: S locally Noetherian, S = B, f flat.
\item Definition artin-definition-openness-versality (canonical tag 07XQ). cited at statement line 1551. The property being established: openness of versality.
\item Lemma artin-lemma-get-openness-obstruction-theory (canonical tag 0CYF). cited at proof lines 1606 and 1630. The openness criterion; its footnote reduces functoriality to fixed A.
\item Lemma lemma-coherent-diagonal (canonical tag 08W6). cited at proof line 1608. Criterion condition (1): representable diagonal.
\item Lemma lemma-coherent-RS-star (canonical tag 08LQ). cited at proof line 1609. Criterion condition (2): Rim-Schlessinger.
\item Lemma lemma-coherent-limits (canonical tag 08KD). cited at proof line 1610. Criterion condition (3): limit preservation.
\item Definition artin-definition-obstruction-theory (canonical tag 07YG). cited at proof line 1626. The obstruction theory axioms.
\item Lemma defos-lemma-flat-ringed-topoi (canonical tag 08VW). cited at proof lines 1624 and 1634. Obstruction class and tangent identification T\_x(M) = Ext\textasciicircum{}1.
\item Lemma defos-lemma-functorial-ringed-topoi (canonical tag 0CYE). cited at proof line 1633. Functoriality of obstruction classes for fixed A.
\item Lemma spaces-perfect-lemma-compute-ext (canonical tag 08JR). cited at proof line 1644. The perfect object K computing Ext\textasciicircum{}i for i <= 2.
\item Lemma more-algebra-lemma-pseudo-coherent-tensor (canonical tag 0CYB). cited at proof line 1662. Product-commutation characterization of pseudo-coherent complexes carries the final step.
\end{itemize}\medskip
\ifdefined\hypertarget\hypertarget{sourcepointer-1}{}\fi
\noindent\textit{Source reference.} Original definition 08KB: sources/stacks-chow/tags.txt, source lines 10881--10881. Cited at statement line 1547: the lemma is phrased in this situation; load-bearing.\par\medskip
\ifdefined\hypertarget\hypertarget{sourcepointer-2}{}\fi
\noindent\textit{Source reference.} Original definition 07XQ: sources/stacks-chow/tags.txt, source lines 10085--10085. Cited at statement line 1551: the property being established; load-bearing.\par\medskip
\ifdefined\hypertarget\hypertarget{sourcepointer-3}{}\fi
\noindent\textit{Source reference.} Original lemma 0CYF: sources/stacks-chow/tags.txt, source lines 16240--16240. Cited at proof lines 1606 and 1630: the criterion applied and the footnote reduction for functoriality; load-bearing twice.\par\medskip
\ifdefined\hypertarget\hypertarget{sourcepointer-4}{}\fi
\noindent\textit{Source reference.} Original lemma 08W6: sources/stacks-chow/tags.txt, source lines 11259--11259. Cited at proof line 1608: condition (1) of the criterion; load-bearing.\par\medskip
\ifdefined\hypertarget\hypertarget{sourcepointer-5}{}\fi
\noindent\textit{Source reference.} Original lemma 08LQ: sources/stacks-chow/tags.txt, source lines 10928--10928. Cited at proof line 1609: condition (2) of the criterion; load-bearing.\par\medskip
\ifdefined\hypertarget\hypertarget{sourcepointer-6}{}\fi
\noindent\textit{Source reference.} Original lemma 08KD: sources/stacks-chow/tags.txt, source lines 10881--10881. Cited at proof line 1610: condition (3) of the criterion; load-bearing.\par\medskip
\ifdefined\hypertarget\hypertarget{sourcepointer-7}{}\fi
\noindent\textit{Source reference.} Original definition 07YG: sources/stacks-chow/tags.txt, source lines 10111--10111. Cited at proof line 1626: the obstruction theory axioms; load-bearing.\par\medskip
\ifdefined\hypertarget\hypertarget{sourcepointer-8}{}\fi
\noindent\textit{Source reference.} Original lemma 08VW: sources/stacks-chow/tags.txt, source lines 11249--11249. Cited at proof lines 1624 and 1634: the obstruction class o(F, F x I, 1) and the tangent identification T\_x(M) = Ext\textasciicircum{}1; load-bearing twice.\par\medskip
\ifdefined\hypertarget\hypertarget{sourcepointer-9}{}\fi
\noindent\textit{Source reference.} Original lemma 0CYE: sources/stacks-chow/tags.txt, source lines 16239--16239. Cited at proof line 1633: condition (ii) functoriality for fixed A; load-bearing.\par\medskip
\ifdefined\hypertarget\hypertarget{sourcepointer-10}{}\fi
\noindent\textit{Source reference.} Original lemma 08JR: sources/stacks-chow/tags.txt, source lines 10859--10859. Cited at proof line 1644: the perfect object K with functorial isomorphisms H\textasciicircum{}i(K x M) = Ext\textasciicircum{}i; load-bearing.\par\medskip
\ifdefined\hypertarget\hypertarget{sourcepointer-11}{}\fi
\noindent\textit{Source reference.} Original lemma 0CYB: sources/stacks-chow/tags.txt, source lines 16236--16236. Cited at proof line 1662: the product-commutation step H\textasciicircum{}i(K x prod M\_n) = prod H\textasciicircum{}i(K x M\_n) is carried by this characterization (K is perfect, hence pseudo-coherent); load-bearing.\par\medskip
\begin{lemma}
\label{lemma-coherent-defo-thy}
In Situation \ref{situation-coherent} assume that
$S$ is a locally Noetherian scheme, $S = B$, and $f : X \to B$ is flat.
Let $\mathcal{X} = \Cohstack_{X/B}$. Then we have openness of
versality for $\mathcal{X}$ (see
Artin's Axioms, Definition \ref{artin-definition-openness-versality}).
\end{lemma}

\begin{proof}[Second proof]
This proof is based on
Artin's Axioms, Lemma \ref{artin-lemma-get-openness-obstruction-theory}.
Conditions (1), (2), and (3) of that lemma correspond to
Lemmas \ref{lemma-coherent-diagonal},
\ref{lemma-coherent-RS-star}, and
\ref{lemma-coherent-limits}.

\medskip\noindent
We have constructed an obstruction theory in the chapter on
deformation theory. Namely, given an $S$-algebra $A$ and an
object $x$ of $\Cohstack_{X/B}$ over $\Spec(A)$ given
by $\mathcal{F}$ on $X_A$ we set
$\mathcal{O}_x(M) = \Ext^2_{X_A}(\mathcal{F}, \mathcal{F} \otimes_A M)$
and if $A' \to A$ is a surjection with kernel $I$, then as obstruction
element we take the element
$$
o_x(A') = o(\mathcal{F}, \mathcal{F} \otimes_A I, 1) \in
\mathcal{O}_x(I) = \Ext^2_{X_A}(\mathcal{F}, \mathcal{F} \otimes_A I)
$$
of Deformation Theory, Lemma \ref{defos-lemma-flat-ringed-topoi}.
All properties of an obstruction theory as defined in
Artin's Axioms, Definition \ref{artin-definition-obstruction-theory}
follow from this lemma except for functoriality of obstruction classes
as formulated in condition (ii) of the definition. But as stated in
the footnote to assumption (4) of
Artin's Axioms, Lemma \ref{artin-lemma-get-openness-obstruction-theory}
it suffices to check functoriality of obstruction classes
for a fixed $A$ which follows from
Deformation Theory, Lemma \ref{defos-lemma-functorial-ringed-topoi}.
Deformation Theory, Lemma \ref{defos-lemma-flat-ringed-topoi}
also tells us that
$T_x(M) = \Ext^1_{X_A}(\mathcal{F}, \mathcal{F} \otimes_A M)$
for any $A$-module $M$.

\medskip\noindent
To finish the proof it suffices to show that
$T_x(\prod M_n) = \prod T_x(M_n)$ and
$\mathcal{O}_x(\prod M_n) = \prod \mathcal{O}_x(M)$.
Apply Derived Categories of Spaces, Lemma
\ref{spaces-perfect-lemma-compute-ext}
to the computation of the Ext groups
$\Ext^i_{X_A}(\mathcal{F}, \mathcal{F} \otimes_A M)$
for $i \leq m$ with $m = 2$. We find a perfect object $K \in D(A)$
and functorial isomorphisms
$$
H^i(K \otimes_A^\mathbf{L} M)
\longrightarrow
\Ext^i_{X_A}(\mathcal{F}, \mathcal{F} \otimes_A M)
$$
for $i = 1, 2$. A straightforward argument shows that
$$
H^i(K \otimes_A^\mathbf{L} \prod M_n) =
\prod H^i(K \otimes_A^\mathbf{L} M_n)
$$
whenever $K$ is a pseudo-coherent object of $D(A)$.
In fact, this property (for all $i$) characterizes
pseudo-coherent complexes, see
More on Algebra, Lemma \ref{more-algebra-lemma-pseudo-coherent-tensor}.
\end{proof}

\end{document}
